
\section{Routing Protocols, Applications, and Challenges}
\label{sec:3}

Routing and information dissemination are fundamental
functions of Vehicular Ad Hoc Networks (VANETs).
Unlike traditional wireless networks, VANETs operate
under rapidly changing connectivity conditions,
high node mobility, and nonuniform vehicle
distributions. These characteristics make the
selection of appropriate routing mechanisms
particularly challenging \cite{Hartenstein2008,Zeadally2012}.

A routing protocol determines how information
is forwarded between communication nodes.
In vehicular environments, a suitable routing
strategy must consider not only network
connectivity but also road topology, vehicle
movement, communication delay, and application
requirements. A protocol that performs well
in a relatively stable network may experience
performance degradation when vehicles move
at high speeds or wireless links become
temporarily unavailable.

Furthermore, not all vehicular applications
require conventional end-to-end routing.
Many safety applications rely on local
broadcasting or geographic information
dissemination, while other services may
benefit from multihop forwarding or
infrastructure-assisted communication.
Consequently, VANET research includes
routing, broadcasting, geographic
forwarding, and delay-tolerant approaches
\cite{Karagiannis2011}.

\subsection{Routing Requirements in VANETs}
\label{sec:3A}

The design of vehicular routing mechanisms
is influenced by several important
network characteristics.

\subsubsection{High Mobility}

Vehicles continuously change their positions
and may move at different speeds and
directions. As a result, the set of
neighboring vehicles can change rapidly.

A communication path that is valid
at one moment may become unavailable
shortly afterward. Routing protocols
must therefore adapt to mobility
without generating excessive
control traffic.

\subsubsection{Dynamic Network Topology}

The network topology depends on
the positions and communication
ranges of participating vehicles.

When vehicles enter or leave
communication range, wireless
links are established or broken.
This may cause route failures,
packet loss, and additional
route discovery operations.

\subsubsection{Variable Traffic Density}

Vehicle density changes according
to road type, traffic conditions,
and time of day.

In sparse networks, communication
may be interrupted because
suitable forwarding vehicles
are unavailable. In dense
networks, excessive transmissions
may increase channel contention
and packet collisions.

\subsubsection{Quality of Service}

Different applications impose
different communication requirements.
For example, safety-related
notifications generally require
timely and reliable delivery,
whereas some information services
can tolerate longer delays.

Routing mechanisms should therefore
consider Quality of Service (QoS)
requirements such as latency,
reliability, throughput,
and communication overhead
\cite{Zeadally2012,Karagiannis2011}.

\subsection{Classification of Routing Protocols}
\label{sec:3B}

VANET routing protocols can be
classified according to the
information used to select
forwarding paths and the
mechanisms employed to
maintain connectivity.

The main categories include
topology-based routing,
position-based routing,
broadcast-based dissemination,
and delay-tolerant forwarding.

\subsubsection{Topology-Based Routing}

Topology-based protocols
use information about
network connections
to establish forwarding paths.

They are commonly divided
into proactive and reactive
routing protocols.

\paragraph{Proactive Routing}

Proactive routing protocols
maintain routing information
continuously, even when
particular routes are not
currently required.

Examples include
Destination-Sequenced
Distance Vector (DSDV)
and Optimized Link State
Routing (OLSR).

DSDV is based on
distance-vector routing
and uses sequence numbers
to help prevent routing loops.
Routing information is
updated as the network changes.

OLSR uses link-state
information and a mechanism
based on Multipoint Relays
(MPRs) to reduce the
overhead associated with
flooding control messages.

The main advantage of
proactive routing is that
routes may already be
available when data
transmission begins.

However, periodic updates
consume wireless resources.
In highly mobile vehicular
environments, frequent
topology changes may
increase control overhead
and reduce route stability.

\paragraph{Reactive Routing}

Reactive protocols establish
routes when communication
is required rather than
maintaining complete routing
information continuously.

Two commonly studied examples
are Ad Hoc On-Demand Distance
Vector (AODV) and Dynamic
Source Routing (DSR).

AODV discovers routes
through route request
and route reply procedures.
When a route is no longer
valid, route maintenance
mechanisms are used to
handle the failure.

DSR uses source routing,
where the sender may
include route information
in the packet header.
Route discovery and
route maintenance are
important components
of its operation.

Reactive routing can
reduce unnecessary
periodic routing traffic.
Nevertheless, route
discovery introduces
additional delay.

Furthermore, frequent
link failures may
cause repeated route
discovery procedures,
which can increase
overhead in high-mobility
environments.

\subsubsection{Position-Based Routing}

Position-based routing
uses geographic location
information to support
packet forwarding.

Instead of relying
entirely on a maintained
end-to-end route, a node
may select a forwarding
neighbor according to
its geographic position
relative to the destination.

Greedy Perimeter Stateless
Routing (GPSR) is a
well-known example
of geographic routing.

In greedy forwarding,
a packet is forwarded
to a neighbor that
provides geographic
progress toward
the destination.

A simplified next-hop
selection rule is

\begin{equation}
n^{*} =
\arg\min_{\substack{
n \in \mathcal{N}(i)\\
d(n,D)<d(i,D)
}}
d(n,D),
\label{eq:greedy}
\end{equation}

where $\mathcal{N}(i)$
represents the neighboring
nodes of vehicle $i$,
$D$ denotes the destination,
and $d(n,D)$ is the
geographic distance between
neighbor $n$ and destination $D$.

The selection is possible
only when at least one
neighbor provides positive
progress toward the destination.
Otherwise, greedy forwarding
reaches a local minimum,
and an alternative recovery
mechanism may be required.

Position-based routing
can reduce dependence
on continuously maintained
global routing tables.
However, it requires
sufficiently accurate
location information.

Its performance may
also be affected by
road geometry, obstacles,
and local connectivity
gaps. Geographic proximity
does not always guarantee
a usable wireless link.

\subsubsection{Broadcast-Based Dissemination}

Broadcast communication
is particularly important
for safety-related VANET
applications.

A vehicle may transmit
a warning message to
multiple nearby vehicles
without establishing
a separate route to
every receiver.

However, uncontrolled
rebroadcasting can
create a broadcast
storm, especially
in dense networks.

Broadcast storms
may increase redundant
transmissions, channel
contention, and
packet collisions.

Possible mitigation
approaches include
probabilistic forwarding,
distance-based forwarding,
and selective rebroadcasting.

\subsubsection{Delay-Tolerant Routing}

Delay-Tolerant Networking
(DTN) approaches are
relevant when continuous
end-to-end connectivity
cannot be guaranteed.

In sparse vehicular
networks, a vehicle
may temporarily store
a packet and forward
it when a suitable
communication opportunity
becomes available.

This strategy is commonly
described as a
store-carry-forward
mechanism.

Although delay-tolerant
forwarding can improve
information delivery
in disconnected networks,
it may introduce
substantial delays.

Therefore, it is
generally unsuitable
for applications that
require immediate
hazard notification,
but it may support
delay-tolerant
information services.

\begin{table*}[t]
\centering
\caption{Qualitative Comparison of Routing Approaches in VANETs}
\label{tab:routing_comparison}
\footnotesize
\begin{tabular}{|p{2.1cm}|p{2.0cm}|p{3.2cm}|p{3.3cm}|p{4.0cm}|}
\hline
\textbf{Category} &
\textbf{Examples} &
\textbf{Main Principle} &
\textbf{Potential Advantage} &
\textbf{Major Limitation} \\
\hline
Proactive &
DSDV, OLSR &
Maintain routing information continuously &
Routes may be available without initial discovery &
Periodic control overhead and stale routes under high mobility \\
\hline
Reactive &
AODV, DSR &
Discover routes when required &
Avoids maintaining all routes continuously &
Route discovery delay and repeated failures \\
\hline
Geographic &
GPSR &
Forward according to geographic position &
Reduced dependence on complete end-to-end route information &
Position errors, local minima, and road obstacles \\
\hline
Broadcast-based &
Selective rebroadcasting &
Disseminate messages to nearby vehicles &
Suitable for local warnings and cooperative awareness &
Redundant transmissions and channel congestion \\
\hline
Delay-tolerant &
Store-carry-forward &
Buffer packets until forwarding becomes possible &
Can operate under intermittent connectivity &
Potentially high delivery delay \\
\hline
\end{tabular}
\end{table*}

\subsection{Routing Performance Considerations}
\label{sec:3C}

Routing performance is influenced
by network topology, mobility,
wireless channel conditions,
and forwarding decisions.

\subsubsection{Route Stability}

A stable communication route
should remain available
for a sufficiently long
period to support
data transmission.

In a simplified
one-dimensional scenario,
the relative speed
between two vehicles
can be represented as

\begin{equation}
v_{\mathrm{rel}} = |v_i-v_j|,
\label{eq:relative_speed}
\end{equation}

where $v_i$ and $v_j$
are the signed velocities
of vehicles $i$ and $j$
along the same road axis.

Under a simplified
constant-speed model,
the time required for
the relative separation
to change by a distance
$\Delta d$ is

\begin{equation}
T = \frac{\Delta d}{v_{\mathrm{rel}}},
\qquad v_{\mathrm{rel}}>0.
\label{eq:contact_time}
\end{equation}

This expression assumes
that the relative motion
changes the separation
monotonically over the
considered interval.

It is an illustrative
mobility relationship,
not a general prediction
of wireless link lifetime.
Actual link duration
also depends on initial
positions, direction
of travel, propagation
conditions, and the
communication range.

\subsubsection{Routing Overhead}

Routing overhead represents
the communication resources
consumed by control messages
and route maintenance.

An illustrative normalized
control overhead metric
may be written as

\begin{equation}
O_{\mathrm{norm}} =
\frac{N_{\mathrm{control}}}
{N_{\mathrm{data,delivered}}},
\label{eq:routing_overhead}
\end{equation}

where $N_{\mathrm{control}}$
is the number of routing
control packet transmissions
and $N_{\mathrm{data,delivered}}$
is the number of successfully
delivered data packets.

This metric must be defined
consistently when comparing
different protocols,
particularly when
retransmissions or
multihop forwarding
are included.

\subsubsection{Routing Trade-Offs}

No single routing approach
is optimal for every
vehicular scenario.

Proactive routing
can reduce initial
route acquisition delay
but introduces periodic
control traffic.

Reactive routing
can avoid maintaining
unused routes, although
route discovery and
link recovery may
increase delay.

Geographic forwarding
can exploit vehicle
location information,
but its performance
depends on positioning
accuracy and local
network connectivity.

Delay-tolerant approaches
can operate in sparse
networks but may not
meet the timing
requirements of
safety applications.

Therefore, routing
protocol selection
should consider
application requirements,
traffic density,
mobility patterns,
and road environments
\cite{Karagiannis2011,Zeadally2012}.

\subsection{Applications of VANETs}
\label{sec:3b}

VANET applications
are commonly classified
into safety-related,
traffic efficiency,
and information services
\cite{Hartenstein2008}.

\subsubsection{Road Safety Applications}

Road safety is a major
motivation for vehicular
communication research.

Potential applications
include emergency
electronic brake warnings,
collision risk alerts,
hazardous road condition
notifications, and
cooperative awareness.

These applications
often involve
time-sensitive messages.
Their effectiveness
depends on timely
delivery, reliable
communication, and
the accuracy of
the exchanged information.

For example, when a
vehicle detects
unexpected braking,
it may generate
a warning intended
for nearby vehicles.

The warning must be
received and processed
within an application-
appropriate time
interval to remain useful.

\subsubsection{Traffic Management Applications}

Traffic management
applications aim to
improve transportation
efficiency through
the exchange of
traffic-related information.

Examples include
congestion monitoring,
dynamic route guidance,
traffic incident
notifications, and
infrastructure-assisted
traffic coordination.

Connected vehicles
may provide information
about local traffic
conditions, while
roadside infrastructure
may distribute
relevant updates.

However, the benefits
of these applications
depend on the quality,
timeliness, and
availability of
the collected information
\cite{Karagiannis2011}.

\subsubsection{Information and Convenience Services}

Non-safety services
may include navigation
assistance, location-based
information, and access
to transportation-related
digital services.

Compared with urgent
safety notifications,
some information services
may tolerate greater
communication delay.

Nevertheless, bandwidth
consumption and
resource allocation
remain important,
particularly when
safety and non-safety
traffic share
communication resources.

\begin{table}[htbp]
\centering
\caption{Application Requirements in VANETs}
\label{tab:vanet_applications}
\footnotesize
\begin{tabular}{|p{1.7cm}|p{2.3cm}|p{2.3cm}|}
\hline
\textbf{Category} &
\textbf{Examples} &
\textbf{Primary Requirement} \\
\hline
Safety &
Emergency braking and hazard alerts &
Timely and reliable delivery \\
\hline
Traffic efficiency &
Congestion information and route guidance &
Information freshness and coverage \\
\hline
Information services &
Navigation and digital services &
Connectivity and service availability \\
\hline
\end{tabular}
\end{table}

\subsection{Security Threats in VANETs}
\label{sec:3D}

Security is a critical
requirement for vehicular
communication because
received information
may influence driving
decisions and traffic
management operations
\cite{Zeadally2012}.

An attacker may attempt
to inject false messages,
disrupt communication,
or obtain sensitive
information about
network participants.

\subsubsection{Sybil Attacks}

In a Sybil attack,
a malicious participant
attempts to represent
multiple identities
within the network.

Such behavior may
distort perceived
traffic conditions
or influence
cooperative decisions.

For example, false
identities could
be used to create
a misleading impression
of vehicle density.

\subsubsection{Denial-of-Service Attacks}

Denial-of-Service
(DoS) attacks attempt
to reduce the
availability of
communication services.

An attacker may
generate excessive
traffic or interfere
with legitimate
message exchange.

Distributed
Denial-of-Service
(DDoS) attacks
involve multiple
sources contributing
to service disruption.

\subsubsection{False Information Injection}

False information
injection occurs
when misleading
or fabricated
messages are
introduced into
the communication
environment.

Examples include
incorrect hazard
reports, manipulated
traffic information,
or false vehicle
status messages.

Even when a message
is associated with
an authenticated
participant, its
content may still
be inaccurate
or malicious.

Therefore, message
authentication alone
does not guarantee
the correctness
of reported events.

\subsubsection{Replay Attacks}

A replay attack
involves retransmitting
previously valid
messages in an
inappropriate context.

For instance,
an old hazard
notification could
be replayed after
the reported
condition no
longer exists.

Freshness checks,
timestamps, and
appropriate message
validation mechanisms
can help reduce
this risk.

\subsubsection{Privacy Threats}

Vehicular messages
may contain information
that reveals
vehicle movement
or location.

Repeated observation
of communication
patterns may create
privacy risks.

Privacy-preserving
mechanisms must
balance participant
accountability with
protection against
unnecessary tracking.

\subsection{Security Requirements and Countermeasures}
\label{sec:3E}

A secure VANET
must address
several complementary
security objectives
\cite{Zeadally2012,Karagiannis2011}.

\textbf{Authentication:}
A receiving entity
should be able to
verify the legitimacy
of relevant message
credentials or
communication participants.

\textbf{Integrity:}
Security mechanisms
should help detect
unauthorized changes
to transmitted messages.

\textbf{Availability:}
Communication services
should remain usable
despite interference,
congestion, and
attempted disruption.

\textbf{Privacy:}
The system should
limit unnecessary
disclosure of
vehicle identity
and movement patterns.

\textbf{Freshness:}
Receivers should
be able to identify
messages that are
outdated or replayed
when the application
requires current
information.

Potential countermeasures
include digital signatures,
certificate-based
authentication,
pseudonym management,
misbehavior detection,
and secure message
validation.

These mechanisms
also introduce
design trade-offs.
For example,
cryptographic
verification consumes
processing resources,
while security-related
message fields
increase communication
overhead.

\begin{table*}[t]
\centering
\caption{Representative VANET Security Threats and Mitigation Approaches}
\label{tab:vanet_security}
\footnotesize
\begin{tabular}{|p{2.6cm}|p{4.0cm}|p{4.0cm}|p{4.0cm}|}
\hline
\textbf{Threat} &
\textbf{Attack Mechanism} &
\textbf{Potential Impact} &
\textbf{Mitigation Approach} \\
\hline
Sybil attack &
Use of multiple deceptive identities &
Manipulated traffic perception or cooperative decisions &
Credential management and misbehavior detection \\
\hline
DoS/DDoS &
Disruption through excessive or malicious traffic &
Reduced communication availability &
Traffic monitoring and resource protection \\
\hline
False information &
Distribution of fabricated vehicle or road data &
Incorrect situational awareness &
Message validation and consistency checks \\
\hline
Replay attack &
Retransmission of outdated messages &
Misleading or stale event information &
Timestamps, freshness checks, and replay protection \\
\hline
Privacy attack &
Observation and correlation of vehicle messages &
Location tracking and identity inference &
Pseudonyms and privacy-aware credential management \\
\hline
\end{tabular}
\end{table*}

\subsection{Major Technical Challenges}
\label{sec:3c}

In addition to routing
and security concerns,
VANETs face several
interconnected
technical challenges
\cite{Hartenstein2008,Karagiannis2011}.

\subsubsection{Communication Reliability}

Wireless communication
is influenced by
signal attenuation,
interference,
obstacles, and
changing connectivity.

Reliable delivery
is particularly
important when
messages are used
to support
safety-related
applications.

\subsubsection{Network Congestion}

High vehicle density
may increase
the number of
simultaneous
transmissions.

Excessive channel
load can lead
to contention,
packet collisions,
and communication
delay.

Congestion control
mechanisms attempt
to regulate
communication
resource usage
while preserving
important information.

\subsubsection{Scalability}

Vehicular networks
must support
different traffic
densities and
deployment sizes.

A mechanism that
performs adequately
in a small scenario
may experience
higher overhead
or reduced
reliability
as the number
of participating
vehicles increases.

\subsubsection{Interoperability}

Connected transportation
systems may contain
vehicles using
different equipment,
communication
technologies,
and protocol
implementations.

Interoperability
is important for
supporting
communication
across heterogeneous
deployments.

\subsubsection{Security and Privacy Trade-Offs}

Strong authentication
and accountability
are important for
detecting malicious
behavior.

However, vehicular
systems must also
protect sensitive
identity and
location information.

Balancing these
requirements remains
an important
research challenge.

\subsection{Discussion}
\label{sec:3F}

The analysis of
routing approaches,
applications, and
security challenges
demonstrates that
VANET design
involves multiple
interdependent
decisions.

Routing protocols
must adapt to
vehicle mobility
and changing
network connectivity.
At the same time,
communication
mechanisms must
satisfy application-
specific requirements
for delay,
reliability,
and coverage.

Security mechanisms
must protect
communication
without introducing
unacceptable
processing or
transmission overhead.

These considerations
suggest that
future VANET
solutions should
be evaluated
under realistic
mobility conditions,
different vehicle
densities, and
representative
wireless environments.

A comprehensive
evaluation should
consider multiple
performance metrics
rather than relying
on a single
indicator of
network quality.
